% =========================================================================
% PATRONES COMUNES EN DOCUMENTOS: LISTAS (CIENCIAS E INGENIERÍA)
% Basado en la Plantilla Académica Sobria y Minimalista.
% =========================================================================
\documentclass[11pt, a4paper]{article}

% --- Paquetes Esenciales ---
\usepackage[utf8]{inputenc}
\usepackage[T1]{fontenc} % Soporte correcto de fuentes vectoriales
\usepackage[spanish, es-tabla]{babel} % Idioma español (es-tabla pone "Tabla" en vez de "Cuadro")
\usepackage{amsmath, amssymb, amsfonts} % Paquete matemático estándar
\usepackage{graphicx} % Para insertar figuras/imágenes
\usepackage{booktabs} % Para tablas con diseño limpio profesional
\usepackage{xcolor} % Soporte de colores
\usepackage{listings} % Para bloques de código de programación

% --- Tipografía de Alta Calidad (Charter) ---
\usepackage{charter}       % Fuente serif Charter: Excelente legibilidad en pantalla e impresión
\usepackage{microtype}     % Ajustes automáticos de espaciado tipográfico (clave del look limpio)
\usepackage{metalogo}      % Corrige el espaciado y kerning de los logos (ej. \LaTeX) con fuentes personalizadas

% --- Diseño de Página y Espaciados ---
\usepackage[margin=2.5cm]{geometry} % Márgenes de 2.5cm equilibrados
\usepackage{setspace}
\setstretch{1.12} % Interlineado relajado para lectura en pantalla sin parecer doble espacio

% --- Paleta de Colores Sóbria ---
\definecolor{primary}{rgb}{0.0, 0.35, 0.35}  % Verde azulado (Teal) apagado - Formal y sobrio
\definecolor{secondary}{rgb}{0.2, 0.2, 0.2} % Gris carbón para texto principal menos agresivo que negro puro
\definecolor{codebg}{rgb}{0.97, 0.97, 0.95}  % Fondo muy suave para fragmentos de código
\definecolor{grayborder}{rgb}{0.85, 0.85, 0.85} % Borde sutil para cajas

% Aplicar color secundario al cuerpo del texto
\makeatletter
\newcommand{\globalcolor}[1]{%
  \color{#1}\global\let\default@color\current@color
}
\makeatother
\AtBeginDocument{\globalcolor{secondary}}

% --- Hipervínculos Elegantes ---
\usepackage{hyperref}
\hypersetup{
    colorlinks=true,
    linkcolor=primary,
    urlcolor=primary,
    citecolor=primary,
    pdftitle={Patrones Comunes en Documentos: Listas},
    pdfauthor={Jair Aguilar Peña}
}

% --- Configuración de Código (Listings) ---
\lstdefinestyle{customcode}{
    backgroundcolor=\color{codebg},
    commentstyle=\color{primary!70}\itshape,
    keywordstyle=\color{primary}\bfseries,
    numberstyle=\tiny\color{gray},
    stringstyle=\color{orange!80!black},
    basicstyle=\ttfamily\small,
    breakatwhitespace=false,
    breaklines=true,
    captionpos=b, % Título del código siempre ABAJO
    keepspaces=true,
    numbers=left,                    
    numbersep=5pt,                  
    showspaces=false,                
    showstringspaces=false,
    showtabs=false,                  
    tabsize=2,
    frame=single,
    rulecolor=\color{grayborder},
    language=[LaTeX]TeX,
    aboveskip=15pt,
    belowskip=10pt
}
\lstset{style=customcode}

% =========================================================================
% DATOS DEL DOCUMENTO
% =========================================================================
\title{\textbf{\color{primary}Patrones comunes en documentos}}
\author{Jair Aguilar Peña}
\date{\today}

% =========================================================================
% INICIO DEL DOCUMENTO
% =========================================================================
\begin{document}

\maketitle

\begin{abstract}
Este documento sirve como una guía de referencia rápida sobre cómo estructurar patrones visuales y organizativos recurrentes en escritos técnicos. En esta sección nos enfocamos en el estudio y modelado de listas (numeradas y no numeradas) siguiendo los principios de legibilidad moderna de la tipografía Charter.
\end{abstract}

\section{Introducción}
Esta sección contiene ejemplos prácticos sobre cómo implementar patrones de presentación visual y estructuración en documentos profesionales de \LaTeX. El propósito de este archivo es servir como una guía de referencia rápida para redactores y estudiantes de ingeniería.

\subsection{Listas}
Las listas son una de las herramientas de documentación más fundamentales para organizar y agrupar información relacionada. Su función principal es facilitar la lectura rápida (escaneo visual) y ordenar conceptos de manera jerárquica.

\subsubsection{Lista Indefinida (No Numerada)}
Una \textbf{lista indefinida} (implementada mediante el entorno \texttt{itemize}) se utiliza cuando el orden de los elementos no es relevante para la comprensión del texto. 

\vspace{0.3cm}
\noindent\textbf{Casos de Uso Comunes:}
\begin{itemize}
    \item Listar los ingredientes necesarios en una receta de cocina (el orden en que los compras o preparas en la mesa no altera el plato).
    \item Enumerar los requisitos previos para cursar una asignatura.
    \item Detallar las características técnicas o componentes de un sistema de hardware.
\end{itemize}

\vspace{0.3cm}
\noindent\textbf{Ejemplo en el Documento:}
\begin{itemize}
    \item[*] Elemento marcado con un asterisco personalizado.
    \item[$\rightarrow$] Elemento marcado con una flecha matemática.
    \item[Nota:] Elemento precedido por una etiqueta de texto personalizada.
\end{itemize} 

La sintaxis para implementar este tipo de lista en \LaTeX\ se detalla a continuación:

\begin{lstlisting}[caption={Código para listas no numeradas personalizadas}]
\begin{itemize}
    \item[*] Elemento marcado con un asterisco personalizado.
    \item[$\rightarrow$] Elemento marcado con una flecha matemática.
    \item[Nota:] Elemento precedido por una etiqueta de texto personalizada.
\end{itemize} 
\end{lstlisting}

\subsubsection{Lista Definida (Numerada)}
Una \textbf{lista definida} (implementada mediante el entorno \texttt{enumerate}) se utiliza cuando el orden de los elementos es crítico y debe seguir una secuencia lógica, temporal o de prioridad específica.

\vspace{0.3cm}
\noindent\textbf{Casos de Uso Comunes:}
\begin{itemize}
    \item Describir los pasos cronológicos para preparar una receta de cocina (p. ej., 1. Mezclar, 2. Hornear).
    \item Detallar los pasos de ejecución de un algoritmo de programación o una demostración matemática.
    \item Estructurar las preguntas de una evaluación o examen para su ponderación secuencial.
\end{itemize}

\vspace{0.3cm}
\noindent\textbf{Ejemplo en el Documento:}
\begin{enumerate}
  \item Paso inicial o concepto prioritario 1.
  \item Paso secuencial o concepto subordinado 2.
\end{enumerate} 

La sintaxis correspondiente para declarar listas numeradas en \LaTeX\ es la siguiente:

\begin{lstlisting}[caption={Código para listas numeradas secuenciales}]
\begin{enumerate}
  \item Paso inicial o concepto prioritario 1.
  \item Paso secuencial o concepto subordinado 2.
\end{enumerate} 
\end{lstlisting}

\end{document}
